\begin{figure}[htbp]
\centering
\begin{tikzpicture}[>=Latex,font=\small,line cap=round]
 \begin{scope}[shift={(0,0)}]
  \node[font=\small\bfseries] at (2,2.65) {Coordinate displacements};
  \draw[->,gray] (-.25,0)--(4.65,0) node[right] {$\Delta v$};
  \draw[->,gray] (0,-.25)--(0,1.9) node[above] {$\Delta\theta^1$};
  \draw[blue!65!black,thick,->] (0,0)--(2,0);
  \draw[blue!65!black,thick] (2,0)--(4,0);
  \draw[blue!65!black,thick,->] (4,0)--(4,.8);
  \draw[blue!65!black,thick] (4,.8)--(4,1.3);
  \draw[blue!65!black,thick,->] (4,1.3)--(2,1.3);
  \draw[blue!65!black,thick] (2,1.3)--(0,1.3);
  \draw[blue!65!black,thick,->] (0,1.3)--(0,.55);
  \draw[blue!65!black,thick] (0,.55)--(0,0);
  \fill (0,0) circle (1.8pt) node[below left] {$x$};
  \node[below] at (2,0) {$l$};
  \node[right] at (4,.65) {$rl$};
 \end{scope}
 \begin{scope}[shift={(7,0)}]
  \node[font=\small\bfseries] at (1.25,2.65) {$F$ components};
  \draw[->,gray] (-.25,0)--(3.25,0) node[right] {$r\Delta v$};
  \draw[->,gray] (0,-.25)--(0,1.9) node[above] {$\Delta\theta^1$};
  \draw[blue!65!black,thick,->] (0,0)--(.9,0);
  \draw[blue!65!black,thick] (.9,0)--(1.5,0);
  \draw[blue!65!black,thick,->] (1.5,0)--(1.5,.9);
  \draw[blue!65!black,thick] (1.5,.9)--(1.5,1.5);
  \draw[blue!65!black,thick,->] (1.5,1.5)--(.6,1.5);
  \draw[blue!65!black,thick] (.6,1.5)--(0,1.5);
  \draw[blue!65!black,thick,->] (0,1.5)--(0,.6);
  \draw[blue!65!black,thick] (0,.6)--(0,0);
  \fill (0,0) circle (1.8pt);
  \node[below] at (.75,0) {$rl$};
  \node[right] at (1.5,.75) {$rl$};
 \end{scope}
\end{tikzpicture}
\caption{A coordinate rectangle for $X=\partial_{\theta^1}$, with $u$
and $\theta^2$ fixed (schematic).  Its sides have equal component scale
in $F$, even though their coordinate scales differ.  The proof averages
translated rectangles over all four base-position variables and the
three bounded momentum labels.}
\label{ho:fig-loops}
\end{figure}
